% =============================================================================
% 0_resumen_y_abstract.tex — Summary and Abstract
% =============================================================================

\chapter*{Resumen}
\addcontentsline{toc}{chapter}{Resumen}
\markboth{Resumen}{Resumen}

Esta memoria documenta el desarrollo y los resultados obtenidos en las sesiones experimentales de la asignatura \Asignatura. El documento abarca el an\'{a}lisis emp\'{i}rico de algoritmos fundamentales, comenzando por la optimizaci\'{o}n mediante descenso por gradiente, donde se explora el impacto de la tasa de aprendizaje en la convergencia, la eficiencia computacional de los mini-batches y la implementaci\'{o}n de criterios de parada. Posteriormente, se estudia el protocolo de Validaci\'{o}n Cruzada (K-Fold) para la selecci\'{o}n de modelos, as\'{i} como las t\'{e}cnicas de regularizaci\'{o}n Ridge (L2) y Lasso (L1) para prevenir el sobreajuste y facilitar la selecci\'{o}n autom\'{a}tica de atributos.

\vspace{0.5cm}
\textbf{Palabras clave:} Aprendizaje Autom\'{a}tico, Descenso por Gradiente, Validaci\'{o}n Cruzada, Regularizaci\'{o}n, Ridge, Lasso.

\newpage

\chapter*{Abstract}
\addcontentsline{toc}{chapter}{Abstract}
\markboth{Abstract}{Abstract}

This report documents the development and results obtained during the experimental sessions of the course \Asignatura. The document covers the empirical analysis of fundamental algorithms, starting with gradient descent optimization, where the impact of the learning rate on convergence, the computational efficiency of mini-batches, and the implementation of tolerance-based stopping criteria are explored. Subsequently, the K-Fold Cross-Validation protocol for model selection is studied, along with Ridge (L2) and Lasso (L1) regularization techniques to prevent overfitting and enable automatic feature selection.

\vspace{0.5cm}
\textbf{Keywords:} Machine Learning, Gradient Descent, Cross-Validation, Regularization, Ridge, Lasso.
